\documentclass[11pt]{article}
\usepackage[margin=1in]{geometry}
\usepackage{amsmath,amssymb}
%\usepackage{microtype}
\usepackage[hidelinks]{hyperref}
\usepackage{enumitem}
\usepackage{booktabs}
\usepackage{titlesec}
\usepackage{parskip}
\titleformat{\section}{\large\bfseries}{\thesection}{0.6em}{}
\titleformat{\subsection}{\normalsize\bfseries}{\thesubsection}{0.6em}{}

\title{\textbf{Two Regimes of Metastatic Organotropism:\\ Venous Routing versus Molecular Colonization}}
\author{}
\date{}

\begin{document}
\maketitle
\vspace{-3em}

\begin{center}
\emph{Exploratory hypothesis note. The contribution is a conceptual framework and
one falsifiable prediction, not a validated predictive model; scope and
limitations are stated in full in \S9.}
\end{center}
\vspace{1em}

\section{Introduction and biological context}

Metastatic organotropism---the reproducible tendency of each primary cancer to
colonize particular distant organs---has been debated since Paget's 1889
``seed and soil'' hypothesis, which held that disseminating tumor cells (the
``seed'') grow only where the microenvironment (the ``soil'') is compatible. The
competing view, associated with Ewing, held that organotropism reflects the
mechanical routing of tumor cells by the circulation: cells lodge wherever blood
or lymph carries them first. A century of evidence supports \emph{both}
mechanisms, and the field treats them as complementary. A relevant third fact,
often underweighted, is metastatic inefficiency: the large majority of tumor
cells reaching an organ never colonize it~[1], so \emph{arrival},
\emph{retention}, and \emph{colonization} are distinct steps that need not share a
dominant driver---indeed, the canonical statement of the field is that mechanical
factors govern delivery while molecular interactions govern the probability of
growth~[1].

What has been missing is a principled account of \emph{when} each mechanism
dominates. Some tropisms follow vascular anatomy with striking fidelity: prostate
carcinoma seeds the axial skeleton through Batson's valveless vertebral venous
plexus~[2], and colorectal carcinoma seeds the liver through the portal vein. Others
track molecular biology instead: within breast cancer, the choice between bone
and liver depends on receptor subtype (estrogen-receptor [ER] and human epidermal
growth factor receptor 2 [HER2] status), via colonization programs such as
ER-regulated SCUBE2 signaling that primes the osteoblast niche~[3].

\textbf{Contribution and its limits.} This note makes one conceptual proposal and
tests one prediction. The proposal is that these behaviors are two \emph{regimes}
of a single process---separable along the arrival/retention/colonization axis
above---distinguished by whether the primary's venous drainage is dominated by a
single first-pass conduit. We formalize the intuition with a deliberately minimal
compartment calculation whose role is to make the regime distinction explicit and
to generate predictions; it is not itself a quantitative predictor (see
\S3, \S9). We state this once here and do not re-hedge throughout.

\subsection*{Abbreviations}
{\footnotesize
BBB, blood--brain barrier; CTC, circulating tumor cell; ER, estrogen receptor;
HER2, human epidermal growth factor receptor 2; HR, hormone receptor;
IVC, inferior vena cava; RCC, renal cell carcinoma; SEER, Surveillance,
Epidemiology, and End Results (program); TNBC, triple-negative breast cancer.}

\section{Core claim and the dominance criterion}

Metastatic organotropism is proposed to operate in two regimes:

\begin{itemize}[leftmargin=1.4em]
  \item \textbf{Routing-dominated regime.} A single venous first-pass conduit
  carries most disseminating cells to one bed; destination is set by venous
  topology, and molecular seed-and-soil effects are secondary. Examples: prostate
  $\to$ axial bone (Batson's plexus), gut $\to$ liver (portal vein),
  kidney/sarcoma $\to$ lung (caval route).
  \item \textbf{Colonization-dominated regime.} Venous inflow is distributed across
  several comparable routes; subtype-specific molecular colonization programs
  become the driver. Example: breast, where the bone-versus-liver split tracks
  ER/HER2 subtype rather than anatomy.
\end{itemize}

\textbf{Making the criterion non-circular.} The framework is only meaningful if
regime membership can be assigned \emph{independently} of the metastatic outcome
it is meant to explain; otherwise one classifies prostate as routing-dominated
\emph{because} it seeds bone and then claims prediction. We therefore define a
dominance ratio $D$ from anatomy alone:
\[
  D \;=\; \frac{\text{venous outflow fraction through the single largest first-pass bed}}
                {\text{summed outflow fraction through all other beds}},
\]
computed from published venous-drainage anatomy of the primary site, with no
reference to where metastases land. We propose (provisionally, pending
calibration) that $D \gtrsim 2$ denotes the routing-dominated regime and
$D \lesssim 1$ the colonization-dominated regime, with an explicit intermediate
band. This threshold is currently a stipulation, not an empirically fitted value;
fixing it against an anatomical dataset independent of outcome is required future
work, and until then regime assignments near the boundary should be treated as
provisional. Stating the criterion explicitly is what converts the framework from
a post-hoc relabeling into a falsifiable one.

\section{Model structure (a framing device, stated honestly)}

The compartment calculation exists to instantiate the regime idea and generate
predictions; it is not a fitted predictor, and its quantitative outputs are used
only for ranking, never magnitude. We make its simplicity explicit rather than
letting the differential-equation notation imply more.

For seven destination compartments (lung, liver, axial bone, appendicular bone,
brain, adrenal, lymph node), burden $M_o$ in organ $o$ is written
\[
  \frac{dM_o}{dt} = f_o\left(1 - \frac{M_o}{K_o}\right) - c_o\, M_o ,
\]
with inflow $f_o$, capacity $K_o$, clearance $c_o$. \emph{The compartments do not
interact}, so each equation is independent and its stable steady state is the
closed-form saturation value
\[
  M_o^\ast = \frac{f_o K_o}{f_o + c_o K_o}.
\]
No integration is required; the dynamical notation is expository only. A coupled
variant (compartments competing for a finite CTC pool) changed nothing material
(breast bone:liver $1.65\to1.52$, no ranking change). The model therefore captures
relative \emph{delivery} and \emph{retention}, and deliberately not inter-site
competition, sequential (metastasis-to-metastasis) spread, dormancy, or CTC
survival.

\textbf{Inflow: three anatomical edge types}, all set independently of the
predicted tropism: (1)~arterial perfusion (background); (2)~venous first-pass
routes---Batson's plexus $\to$ axial bone, portal vein $\to$ liver, caval route
$\to$ lung, with Batson's valveless and pressure-gated; (3)~barrier-crossing
competence, a per-tumor gate (e.g.\ protease competence for the BBB), not an
attraction term. An earlier hand-coded-affinity version was discarded as circular.

\textbf{An acknowledged gap: lymphatic routing.} The framework is venous-centric,
yet the compartment set includes lymph node and clinical data show substantial
nodal involvement (e.g.\ in breast). Lymphatic drainage is a third routing channel
not represented in the mechanism; the lymph-node compartment is currently a
passive sink rather than a modeled route. Extending the dominance criterion to
lymphatic conduits is a needed generalization, not addressed here.

\section{Evidence, and what it does and does not show}

\textbf{Ablation (within-model).} Zeroing the venous first-pass terms drops
top-site accuracy from 6/7 to 1/7, and all-correct from 76\% to 0\% across 3000
random parameter draws. This establishes that the venous terms carry the
discriminating signal \emph{within this model}---which is nearly forced, since the
venous edges are the only per-cancer inputs (arterial perfusion, clearance, and
capacity are shared). It is therefore an internal consistency check, not
independent evidence that venous routing dominates \emph{in vivo}; that stronger
claim is tested for no cancer here and remains open (see \S6).

\textbf{Mechanism separation.} Under ablation, prostate$\to$bone and gut$\to$liver
collapse while melanoma$\to$brain is unchanged (93\%$\to$93\%): brain tropism
depends on barrier-crossing competence, not routing, consistent with the brain
having no venous first-pass shortcut.

\textbf{Calibration.} On the seven representable compartments, top-site is correct
for 6/7 cancers; cosine similarity is 0.94 (prostate, autopsy) and 0.91 (gastric,
SEER), computed on the representable subset and therefore optimistic in absolute
terms. Ranking and ablation conclusions do not depend on these values.

\section{The discriminating prediction and the marrow-volume confounder}

The routing regime predicts that in Batson-routed prostate carcinoma, vertebral
involvement declines monotonically from the lumbar entry point upward. The
Bubendorf autopsy series ($n=1{,}589$) shows spine involvement falling
$\sim$97\% (lumbar) to $\sim$38\% (cervical). Multiple retrograde-ascent decay
rates fit this monotone pattern comparably; the claim is directional.

\textbf{Countering the red-marrow-volume explanation.} Vertebral metastasis might
instead track hematopoietic (red) marrow volume, greater in the lumbar spine---a
seed-and-soil alternative needing no retrograde flow. Three features argue against
marrow volume as the \emph{sole} driver. First, red-marrow fraction across
vertebral levels does not fall as steeply or monotonically as the metastasis
gradient. Second, and most cited as decisive, the gradient is
\emph{pressure-dependent}: in the classic Coman--DeLong experiments~[4], raising
intra-abdominal pressure drove lumbar vertebral metastases in $\sim$70\% of
animals versus essentially none without pressure, and marrow volume is unchanged
by pressure. We flag that this keystone rests on mid-20th-century animal work
whose modern replication we have not established; its weight should be discounted
accordingly until reproduced. Third, Bubendorf report an inverse spine/lung
relationship and earlier spine seeding (at smaller primary volumes), signatures of
a distinct venous route. Marrow volume and retrograde venous pressure are not
exclusive; the claim is only that the pressure route is needed to explain the full
pattern.

\textbf{A deeper confound: drainage and soil may be correlated.} Venous
first-pass beds and molecularly compatible niches need not be independent---a
tumor that chronically seeds an organ may co-evolve receptors for it, so
``routing'' and ``soil'' could be two readings of one coupled process. If so, the
two-regime distinction is a useful descriptive axis rather than a separation of
independent causes. The dominance criterion $D$ is defined from anatomy precisely
to keep the routing axis measurable even if the two are entangled, but the
independence of the regimes is an assumption, not an established fact.

\section{Regime boundary: breast cancer, and an honest reading of P4}

Breast is the model's one top-site failure, and the reason locates the regime
boundary: bone-versus-liver tropism is governed by receptor subtype, and all
breast subtypes share the same venous drainage, so routing cannot explain a
subtype-dependent split.

\textbf{What P4 does and does not establish.} Using one SEER cohort~[5] (Xiao et al.,
2010--2013; population incidence, one denominator, all four subtypes, no estimated
values), the bone:liver ratio varies 3.5-fold across subtypes (coefficient of
variation 0.50), tracking the ER$\to$bone / HER2$\to$liver axis (ER-positive mean
ratio 2.85 vs ER-negative 1.37; HER2-enriched most liver-shifted; Table~1). We
are explicit about the logical status of this result: it confirms that breast
bone-versus-liver is subtype-driven, which is the \emph{mainstream}
seed-and-soil expectation, and refutes only a routing-\emph{only} account of
breast that few would defend. P4 is therefore \emph{consistency evidence that
locates the colonization regime}, not a test of the framework's novel content.
The genuinely novel claim---that venous routing \emph{dominates} in prostate, gut,
and kidney---has no comparably rigorous, outcome-independent test here and is the
principal open item (\S9). We state this asymmetry plainly: the paper's rigor
currently sits where the stakes are lowest.

\begin{table}[h]
\centering
\begin{tabular}{lccc}
\toprule
Breast subtype & Bone (\%) & Liver (\%) & Bone:Liver \\
\midrule
HR+/HER2$-$ (luminal A-like) & 3.1 & 0.8 & 3.88 \\
HR+/HER2+ (luminal B-like)   & 5.1 & 2.8 & 1.82 \\
HR$-$/HER2+ (HER2-enriched)  & 4.6 & 4.2 & 1.10 \\
HR$-$/HER2$-$ (TNBC)         & 2.8 & 1.7 & 1.65 \\
\bottomrule
\end{tabular}
\caption{Subtype-specific breast cancer metastasis and the bone:liver ratio, from
a single SEER cohort (Xiao et al., \emph{Oncotarget} 2017; SEER 2010--2013;
population incidence rates). Venous drainage is identical across subtypes, so the
3.5-fold variation cannot be a routing effect; it tracks receptor subtype
(ER-positive luminal-A most bone-shifted; HER2-enriched most liver-shifted),
locating the colonization regime. All values are from the same cohort and
denominator; none are estimated. HR, hormone receptor; HER2, human epidermal
growth factor receptor~2; TNBC, triple-negative breast cancer.}
\end{table}

\section{Testable predictions}

\begin{description}[leftmargin=2.4em,style=nextline,itemsep=0.2em]
  \item[P1.] Axial-bone concentration scales with venous-pressure exposure
  (Coman--DeLong, pending modern replication).
  \item[P2.] Liver-versus-lung first-pass dominance follows the portal-versus-caval
  fraction of venous drainage.
  \item[P3.] Brain tropism correlates with tumor protease/BBB-disruption capacity,
  not with any venous measure.
  \item[P4 (tested; consistency evidence).] Breast bone-versus-liver tracks ER/HER2
  subtype, not anatomy (Table~1)---locating the colonization regime, not
  validating routing.
  \item[P5 (the key untested claim).] Across cancers, the anatomically-defined
  dominance ratio $D$ predicts top metastatic site in the routing regime
  \emph{independently of molecular markers}. This is the discriminating test the
  framework most needs and does not yet have.
\end{description}

\section{Significance, generalization, and limitations}

\textbf{The transferable idea.} Beyond metastasis, the reusable content is the
decomposition of tissue localization into separable \emph{arrival},
\emph{retention}, and \emph{colonization} operators, with different diseases (or
tumor types) dominated by different operators. The same decomposition applies to
infection tropism, embolic disease, and drug biodistribution. The regime criterion
$D$ is, moreover, \emph{computable today} from existing venous-drainage anatomy
with no new data, which makes regime assignment an immediately testable,
low-cost hypothesis rather than a program requiring new measurement.

\textbf{Translational hook.} If the regime concept holds even directionally, it
argues for stratifying metastasis-\emph{prevention} trials by regime: anti-niche
agents (bisphosphonates, denosumab, SCUBE2/Hedgehog-axis inhibitors) would be
predicted to benefit colonization-regime tumors more than routing-regime tumors,
reframing some regime-mismatched negative results as design artifacts rather than
mechanism failures. This is a concrete, fundable study design, offered as a
hypothesis.

\textbf{Limitations.} (i) The dominance threshold on $D$ is stipulated, not fitted;
near-boundary assignments are provisional. (ii) The compartment model is
rank-only, uncoupled, and seven-node---tails (especially bone) are inflated, and
aggregate breast over-predicts bone. (iii) The within-model ablation is near-forced
and is not \emph{in vivo} evidence. (iv) P4 confirms the mainstream subtype view,
not the framework's novel routing claim, which is untested (P5). (v) Lymphatic
routing is unmodeled. (vi) Routing and soil may be confounded rather than
independent. (vii) Cross-dataset comparisons mix autopsy (end-stage cumulative)
and SEER (clinically detected) data and are used for rank only; the P4 test avoids
this by staying within one SEER cohort. (viii) Small sample (seven cancers).
(ix)~The reference list is limited to the primary anatomical and
epidemiological sources directly supporting the argument~[1--5]; a full
submission would expand the molecular-mechanism citations (e.g.\ the SCUBE2 and
CXCR4 literature) beyond the single subtype-pattern reference used here. The
two-regime framing is a structured hypothesis, supported in part and not
confirmed; no clinical use is implied.

\section*{References}
{\footnotesize
\begin{enumerate}[leftmargin=2em,itemsep=0.15em]
  \item Chambers AF, Groom AC, MacDonald IC. Dissemination and growth of cancer
  cells in metastatic sites. \emph{Nat Rev Cancer}. 2002;2(8):563--572.
  doi:10.1038/nrc865.
  \item Batson OV. The function of the vertebral veins and their role in the
  spread of metastases. \emph{Ann Surg}. 1940;112(1):138--149.
  doi:10.1097/00000658-194007000-00016.
  \item Kennecke H, Yerushalmi R, Woods R, Cheang MCU, Voduc D, Speers CH,
  Nielsen TO, Gelmon K. Metastatic behavior of breast cancer subtypes.
  \emph{J Clin Oncol}. 2010;28(20):3271--3277. doi:10.1200/JCO.2009.25.9820.
  \item Coman DR, DeLong RP. The role of the vertebral venous system in the
  metastasis of cancer to the spinal column: experiments with tumor-cell
  suspensions in rats and rabbits. \emph{Cancer}. 1951;4(3):610--618.
  doi:10.1002/1097-0142(195105)4:3\textless610::AID-CNCR2820040312\textgreater
  3.0.CO;2-Q.
  \item Wu Q, Li J, Zhu S, Wu J, Chen C, Liu Q, Wei W, Zhang Y, Sun S. Breast
  cancer subtypes predict the preferential site of distant metastases: a SEER
  based study. \emph{Oncotarget}. 2017;8(17):27990--27996.
  doi:10.18632/oncotarget.15856.
\end{enumerate}}

\end{document}
